\section{A sharp theorem for every two-tree}\label{sec:twotree}
A $k$-tree starts with $K_{k+1}$ and repeatedly adds a vertex whose neighbors
are the vertices of an existing $k$-clique. In particular, a two-tree is built
by adding triangle ears on edges. This section has no diameter restriction.

\begin{theorem}[Exact minimum boundary of a two-tree]\label{thm:twotree}
Among all two-trees on $n$ vertices, the minimum boundary cardinality is
\[
 \min b(G)=
 \begin{cases}
 3,&n=3,\\
 4,&n=4,5,\\
 \lceil n/2\rceil+2,&n\ge6.
 \end{cases}
\]
In particular, $n\le2b(G)-4$ for every two-tree of order at least six.
There are exactly three two-trees with five boundary vertices: the book of
three triangles $K_2\vee\overline{K_3}$, the cone over $P_5$, and the cone
over the five-vertex tree of degree sequence $(3,2,1,1,1)$.
\end{theorem}
The order bound is sharp at every even order, and its rounded boundary form
is sharp at every odd order as well. Maximality matters: there are
two-connected series-parallel graphs on nine vertices with boundary six,
so the same order bound does not hold for all partial two-trees.

\subsection{Localizing stability in chordal graphs}
A graph is chordal if it has no induced cycle of length at least four.
We use the elementary chordal facts that minimal vertex separators are
cliques and that a perfect elimination ordering can end in any specified
clique.

\begin{theorem}[Radius-two localization]\label{thm:local}
In every connected chordal graph,
\[
 \stab_G(u)=\max_{y:d_G(u,y)\le2}\stab_G(u,y).
\]
This identity concerns the maximum stability itself, not just its sign.
\end{theorem}
\begin{proof}
Let $y\notin N[u]$, let $C$ be its component in $G-N[u]$, and put
$S=N(C)\subseteq N(u)$. For every $s\in S$ a path $u,s,C,y$ avoids
$S\setminus\{s\}$, so $S$ is a minimal $u$--$y$ separator and is a clique.
There is a vertex $x\in C$ adjacent to all of $S$. For completeness, take
a perfect elimination ordering of $G[S\cup C]$ ending in $S$.
Deleting a simplicial vertex of $C$ preserves connectivity of the remaining
$C$, and preserves its full neighborhood $S$: a neighbor $s\in S$ of the
deleted vertex is also adjacent to its remaining neighbor in $C$.
The last vertex of $C$ is therefore complete to $S$.

Set $r=\min_{s\in S}d(s,y)$ and $A=\{s\in S:d(s,y)=r\}$.
Since $S$ is a clique, its other members have distance $r+1$.
Every path from $N(u)$ to $y$ enters $C$ through $S$. It follows that the
distances of vertices of $N(u)$ to $y$ are $r$ on $A$, $r+1$ on
$(N(u)\cap N[A])\setminus A$, and $r+2$ on $N(u)\setminus N[A]$.
As $d(u,y)=r+1$,
\[
 \stab(u,y)=|A|-|N(u)\setminus N[A]|
 \le |S|-|N(u)\setminus N[S]|=\stab(u,x).
\]
The final equality uses $N(x)\cap N(u)=S$, and $d(u,x)=2$.
\end{proof}

\begin{lemma}[Interior degree in a chordal graph]\label{lem:chordal-degree}
In a $k$-connected chordal graph, $k\ge2$, every nonuniversal interior
vertex has degree at least $2k+1$.
\end{lemma}
\begin{proof}
First $G[N(u)]$ is connected. A path in $G-u$ between two neighbors of $u$
can be shortened into $N(u)$, because each excursion through a component
of $G-N[u]$ enters and exits its clique separator.
Choose such a component, with separator $S$ of size $s\ge k$, and the
vertex $x$ complete to $S$ from the preceding proof. Put
$T=N(u)\cap N[S]$. Interiorness gives
$0\ge\stab(u,x)=s-\dd(u)+|T|$, so $\dd(u)\ge s+|T|$.
The set $S$ is proper in $N(u)$, since otherwise that score is positive.
Connectivity of $G[N(u)]$ gives $|T|\ge s+1$, proving the result.
\end{proof}

For an edge $ab$ of a two-tree, put $t_{ab}=|N(a)\cap N(b)|$.
The neighbor graph $G[N(u)]$ is a tree; the degree of its vertex $v$ is
$t_{uv}$. This follows directly by induction under triangle-ear addition.
Each exterior component has a separator which is an edge $ab$ of this tree;
such an attachment exists exactly when $t_{ab}\ge2$. Consequently
Theorem~\ref{thm:local} gives the exact formula
\begin{equation}\label{eq:local-two}
\stab(u)=\max\left\{
\begin{array}{ll}
2+t_{uv}-\dd(u)&(v\in N(u)),\\
2+t_{ua}+t_{ub}-\dd(u)&(uab\text{ a triangle},\ t_{ab}\ge2).
\end{array}\right.
\end{equation}
Indeed, the closed neighborhood of the edge $ab$ in the neighbor tree has
size $t_{ua}+t_{ub}$. A common exterior neighbor of $a,b$ gives the second
score; a neighboring target $v$ gives the first.

\subsection{Pruning and the sharp count}
\begin{lemma}[Pruning]\label{lem:pruning}
Let $G$ be a noncomplete $k$-tree, $k\ge2$, and remove any set $Q$ of
degree-$k$ vertices. If $H=G-Q$ has at least $k+2$ vertices, every
simplicial vertex of $H$ is a boundary vertex of $G$.
\end{lemma}
\begin{proof}
Degree-$k$ vertices of a noncomplete $k$-tree are independent and simplicial.
Their removal leaves a $k$-tree while at least $k+1$ vertices remain.
For a simplicial $u$ of $H$, the $k$-clique $S=N_H(u)$ extends in $H-u$
to a clique $S\cup\{x\}$. Thus $x\notin N_H[u]$ and $d_G(u,x)=2$.
Every removed neighbor $q$ of $u$ has its other $k-1$ neighbors in $S$.
Since $k\ge2$, $d_G(q,x)=2$. The $k$ vertices of $S$ are closer and
all other neighbors are equidistant, so $\stab_G(u,x)=k>0$.
\end{proof}
If the remainder is $K_k$, all its vertices are universal. If it is
$K_{k+1}$ and all its vertices were interior, a removed vertex $q$ omits
one core vertex $u$ and gives the same positive score $\stab_G(u,q)=k$.
These observations handle the small remainders below.

\begin{proof}[Proof of the lower bound in Theorem~\ref{thm:twotree}]
Suppose $n\ge6$, let $I=V\setminus\bd G$, let $i=|I|$, and write
$I_5=\{u\in I:\dd(u)=5\}$. Every interior vertex has degree at least
five, by Lemma~\ref{lem:chordal-degree}, or by $\dd(u)=n-1\ge5$
when it is universal.

We assign a distinct boundary vertex to every $u\in I_5$. Its neighbor
tree has five vertices. It cannot be a star, because the star center gives
a positive first score in \eqref{eq:local-two}. If the neighbor tree is
$P_5$, the second score says that only its two terminal edges can have
exterior attachments: an active edge's endpoint degrees must sum to at
most three. Its middle vertex $c$ therefore has global degree three.
If the neighbor tree is the fork of degrees $(3,2,1,1,1)$, only the edge
from the degree-two vertex to its leaf can have an exterior attachment.
The branching vertex $c$ has global degree four. In both cases $c$ is
boundary by Lemma~\ref{lem:chordal-degree}; assign it to $u$.

This assignment is injective. The two types cannot collide because their
assigned vertices have different global degrees. In the path case, write
the neighbor path as $a,b,c,d,e$. The only alternative interior neighbors
that could assign $c$ are $b,d$. In either of their neighbor trees, $c$
has degree one, rather than the degree two required of a path middle.
In the fork case, two other neighbors of the assigned branch have global
degree two. The sole remaining possible interior neighbor again sees
the assigned vertex with link degree one, rather than three.

Set $E_B=\sum_{v\in\bd G}(\dd(v)-2)$. Injectivity gives
$E_B\ge|I_5|$: path assignments use one unit, and fork assignments two.
The degree sum $4n-6$ now yields
\begin{equation}\label{eq:charge}
2(b-i-3)=E_B-|I_5|+
 \sum_{u\in I:\dd(u)\ge7}(\dd(u)-6)\ge0.
\end{equation}
Equality would force all assignments to be path assignments, and every
boundary vertex of degree greater than two to be an assigned degree-three
middle. Each such middle has three neighbors of degree greater than two:
its assigned interior root and its two neighboring vertices in the path.

Delete all degree-two vertices. They are independent, and the remainder
$H$ has at least two vertices, because each deleted vertex has two
undeleted neighbors. Every remaining boundary vertex still has degree
three in $H$. If $H$ is a noncomplete two-tree, its simplicial degree-two
vertex cannot be boundary, but cannot be interior by Lemma~\ref{lem:pruning}.
If $H=K_3$, all its vertices would be interior, excluded by the small-remainder
argument. If $H=K_2$, both vertices are universal and boundary, again
impossible. Equality in \eqref{eq:charge} is therefore excluded.
Its integer right side is even, hence at least two. Thus $b-i\ge4$,
and $n=b+i\le2b-4$.
\end{proof}

\subsection{Matching constructions and classification at five}
For $r\ge2$, start with a path $v_1,\ldots,v_r$ and independent vertices
$c_1,\ldots,c_r$. Join $c_j$ to $v_j$ and its available path neighbors
$v_{j-1},v_{j+1}$. This core is a two-tree. At each endpoint $v$, with
its corresponding $c$, choose either of the following three-vertex-end
configurations (the vertex $c$ is already present):
\begin{itemize}
\item add $p$ on $vc$, and then $q$ on $vp$;
\item add $p,q$ separately on $vc$.
\end{itemize}
Every $v_j$ has degree five; its neighbor tree is a path or the permitted
fork. Formula~\eqref{eq:local-two} gives $\stab(v_j)=0$.
Every other vertex has degree at most four and is boundary. Thus
$n=2r+4$ and $b=r+4$. The three unordered choices of endpoint types
give three equality examples for each $r\ge2$. The following sections prove
that these exhaust the even-order equality cases, classify all odd-order
cases, and give their exact unlabelled counts.
For $r=1$ the two appropriate choices give the cones over $P_5$ and
the fork, of order six. To obtain every odd order, use the first type at
both ends and add one ear on the outer edge $pq$. This activates only
a permitted terminal edge in the adjacent interior vertex's link.
The old $p,q$ have stability at least two, so remain boundary, and exactly
one new boundary vertex is added. This proves sharpness at every order.

At boundary five the bound forces $n\le6$. The two two-trees of order
five are the three-triangle book and the cone over $P_4$, of boundaries
five and four respectively. At order six there is one interior vertex;
its degree is five, so it is universal. A two-tree with a universal vertex
is the cone over a tree. The only five-vertex trees without a universal
vertex are $P_5$ and the fork. This proves the stated classification.

\begin{proposition}[All $k$-trees at boundary five]\label{prop:ktrees}
Every $k$-tree, $k\ge2$, satisfies
\[
 n\le(k+1)(b-k),
\]
with equality only for $K_{k+1}$. Across all $k\ge2$, the five-boundary
graphs are the three two-trees above, $K_5$ minus one edge (a three-tree),
and $K_5$ (a four-tree). For $k=1$ they are all trees with five leaves.
\end{proposition}
\begin{proof}
When $n\ge2k+2$, interior degrees are at least $2k+1$ and boundary
degrees at least $k$. The degree sum $2kn-k(k+1)$ gives the bound.
There are at least two boundary vertices of degree greater than $k$:
delete all degree-$k$ vertices and apply Lemma~\ref{lem:pruning} to the
two simplicial vertices of a noncomplete remainder. For a complete
remainder, at most one boundary core vertex would force every deleted
vertex to omit that same vertex, leaving at least two universal interior
vertices, a contradiction. Hence the large-order bound improves by two.

For $n\le2k+1$, minimum degree $k$ forces diameter at most two.
The known diameter-two theorem gives $b\ge n-1$. The cases
$n=k+1,k+2$ are a clique and a clique minus one edge; in the latter
there are at least two universal vertices and $b=n$.
For $n\ge k+3$,
$(k+1)(n-k-1)-n=kn-(k+1)^2\ge k-1>0$.
This proves the bound and its equality statement.
For $b=5$, it leaves $k=2,3,4$ and orders at most $9,8,5$.
Independent ear generation and recursive simplicial-degree-$k$
recognition on those finite ranges give exactly the five listed graphs.
The two-tree part has already been proved directly above.
\end{proof}

